\labsection{7}{Генетический алгоритм с перестановочным представлением особей. Задача коммивояжёра}

\labpart{Цель работы}
Изучить представление особей в форме перестановок, специализированные операторы скрещивания и мутации
для него, особенности решения задач комбинаторной оптимизации генетическими алгоритмами и способы учёта
ограничений; реализовать ГА, решающий задачу коммивояжёра, и меметический алгоритм с локальным поиском.

\labpart{Теоретические сведения}
Раздел~\ref{sec:combinatorial} теоретической части: учёт ограничений в ГА на примере задач о рюкзаке
(\ref{eq:knapsack}) и о минимальном покрытии (\ref{eq:cover}), задача коммивояжёра (\ref{eq:tsp}),
представления тура, операторы PMX, OX, CX, мутации перестановок, локальный поиск 2-opt (\ref{eq:2opt}),
меметические алгоритмы.

\labpart{Постановка задачи}
Дано множество городов, заданных координатами $(x_i; y_i)$ на плоскости; стоимость
перехода между городами равна евклидову расстоянию. Построить генетическим алгоритмом тур обхода всех
городов минимальной длины, используя представление путей, оператор кроссинговера и тестовый пример из
таблицы~\ref{tab:l7-var}. Сравнить найденный тур с оптимальным из файла \code{*.opt.tour} и
графически проиллюстрировать поиск.

\begin{table}[!htb]
\caption{Варианты заданий к лабораторной работе №7}
\label{tab:l7-var}
\centering\small
\begin{tabular}{clccc}
\toprule
№ & Оператор кроссинговера & Файл TSPLIB & Городов & Оптимальный тур \\
\midrule
1 & частично соответствующий (PMX) & \code{berlin52.tsp} & 52  & 7542 \\
2 & упорядоченный (OX)             & \code{eil51.tsp}    & 51  & 426 \\
3 & циклический (CX)               & \code{eil76.tsp}    & 76  & 538 \\
4 & частично соответствующий (PMX) & \code{pr76.tsp}     & 76  & 108\,159 \\
5 & упорядоченный (OX)             & \code{st70.tsp}     & 70  & 675 \\
6 & циклический (CX)               & \code{lin105.tsp}   & 105 & 14\,379 \\
7 & частично соответствующий (PMX) & \code{kroA100.tsp}  & 100 & 21\,282 \\
\bottomrule
\end{tabular}
\end{table}

Файлы TSPLIB~\cite{reinelt1991} текстовые: заголовок (\code{NAME}, \code{TYPE : TSP}, \code{DIMENSION}~---
число городов, \code{EDGE\_WEIGHT\_TYPE : EUC\_2D}~--- евклидово расстояние, округлённое до целого),
раздел \code{NODE\_COORD\_SECTION} со строками <<номер $x$ $y$>> и строка \code{EOF}. В файле
\code{*.opt.tour} в разделе \code{TOUR\_SECTION} перечислены города оптимального тура, список
завершается числом $-1$; пересчёт его длины программой проверяет чтение данных и функцию расстояния.

\labpart{Требования к программе}
Исходные данные: размер популяции, вероятности кроссинговера и мутации, максимальное число поколений,
файл \code{*.tsp} со списком городов. Отображаются лучший тур текущей популяции в виде графа на
плоскости и в виде списка городов, графики лучшей и средней длины тура по поколениям и длина лучшего
тура. Отладку проводить на файлах TSPLIB, сравнивая решения с оптимальными
турами.

\labpart{Порядок выполнения работы}
\begin{steps}
\item Реализовать чтение файлов \code{*.tsp} и \code{*.opt.tour}, матрицу расстояний с
округлением по правилу \code{EUC\_2D} и длину тура (подраздел~\ref{sec:py-tsp}); проверить, что длина опубликованного оптимального тура
совпадает с указанной в таблице~\ref{tab:l7-var}.
\item Реализовать ГА с представлением путей: оператор кроссинговера варианта, мутация обменом или
инверсией, турнирная селекция, элитарное сокращение без дубликатов. Проверить оператор кроссинговера на
примере из подраздела~\ref{sec:tsp} (для PMX и OX~--- также сравнением с DEAP).
\item Выполнить серию запусков и сравнить лучшую найденную длину с оптимальной.
\item Дополнительно: добавить локальный поиск 2-opt к каждому потомку (меметический ГА) и сравнить с
мультистартом 2-opt без ГА при одинаковом времени работы.
\end{steps}

\labpart{Пример выполнения}
Вариант~3: набор eil76 (76 городов, оптимальный тур длиной 538, средняя длина
случайного тура~--- 2520), циклический кроссинговер. Результаты приведены в таблице~\ref{tab:l7-tsp} и на
рисунке~\ref{fig:l7-tours}.

\begin{table}[!htb]
\caption{Задача коммивояжёра eil76, оптимум 538 (ГА~--- $150 \times 500$, меметический~--- $60 \times 60$)}
\label{tab:l7-tsp}
\centering\small
\setlength{\tabcolsep}{4pt}
\begin{tabular}{lccccc}
\toprule
Метод & Запусков & Лучший & Медиана & Худший & Отклонение, \% \\
\midrule
ГА: CX + обмен & 10 & 772 & 818 & 847 & $+52,0$ \\
ГА: CX + инверсия & 10 & 876 & 965 & 1071 & $+79,4$ \\
Меметический ГА & 5 & 538 & 538 & 538 & 0 \\
Мультистарт 2-opt & 5 & 540 & 550 & 553 & $+2,2$ \\
\bottomrule
\end{tabular}
\end{table}

\begin{figure}[!htb]
\centering
\includegraphics[width=\textwidth]{\labfig{7}{B1_tours.png}}
\caption{Лучшие найденные маршруты и опубликованный оптимальный тур eil76}
\label{fig:l7-tours}
\end{figure}

ГА с CX за 75\,000 оценок даёт туры на 43--57\,\% длиннее оптимального: CX сохраняет позиции городов,
а для маршрута важна их смежность, и рёбра хороших родителей почти не наследуются. Меметический ГА нашёл тур длиной 538 во всех пяти запусках
примерно за 5~с (маршрут отличается от опубликованного: у eil76 несколько оптимальных туров), а
мультистарт 2-opt за то же время~--- лишь 540--553.

\labpart{Контрольные вопросы}
\begin{questions}
\item Сформулируйте задачи о рюкзаке, о минимальном покрытии и о коммивояжёре. Какое представление
решения используется в каждой из них?
\item Какие способы учёта ограничений применяются в ГА? Сравните штрафы, восстановление и декодеры на
примере задачи о рюкзаке.
\item Приведите способы представления тура и их особенности. Почему для задачи коммивояжёра не подходит
обычный двоичный кроссинговер?
\item Опишите представление соседства и операторы обмена рёбер, обмена подтуров и эвристического
скрещивания.
\item Опишите порядковое представление тура. Какой оператор кроссинговера для него применим?
\item Опишите операторы PMX, OX и CX для представления путей; выполните их для родителей
$(1\ 2\ 3\ 4\ 5\ 6\ 7\ 8)$ и $(3\ 7\ 5\ 1\ 6\ 8\ 2\ 4)$.
\item Какие мутации применяются к перестановкам? Опишите алгоритм 2-opt. Что такое меметический алгоритм и почему он эффективнее чистого ГА?
\end{questions}
